\documentclass[12pt,letterpaper]{article}
\setlength{\parindent}{0pt}
\setlength{\parskip}{1ex plus 0.5ex minus 0.2ex}
\addtolength{\hoffset}{-1in}
\addtolength{\textwidth}{2in}
\addtolength{\voffset}{-1.3in}
\addtolength{\textheight}{2in}

\usepackage{amsmath,amssymb,amsthm}
\usepackage{booktabs}

\begin{document}
\par
{\bfseries Behrouz Barati B} \\
https://github.com/1ncompleteness/Numerical-Analysis
\par
\begin{center}
\bfseries Homework 5 \\
\bfseries Math.~481a, Spring 2026 \\
\bfseries Prepared in \LaTeX\
\end{center}

%% PROBLEM 1

\textbf{Problem 1.}

\textit{Claim.} If $f''(x_0)$ exists and $f$ is defined in some open interval containing $x_0$, then
\[
\lim_{h \to 0} \frac{1}{h^2}\left[f(x_0 - h) - 2f(x_0) + f(x_0 + h)\right] = f''(x_0).
\]

\textit{Proof.} Since $f''(x_0)$ exists, $f'$ exists on some open interval $(x_0 - \delta, x_0 + \delta)$, and by definition of the second derivative,
\[
\lim_{h \to 0} \frac{f'(x_0 + h) - f'(x_0)}{h} = f''(x_0), \qquad \lim_{h \to 0} \frac{f'(x_0 - h) - f'(x_0)}{-h} = f''(x_0).
\]

Define $\varphi(h) = f(x_0 + h) - 2 f(x_0) + f(x_0 - h)$. Then $\varphi(0) = 0$, and for $|h| < \delta$, $\varphi$ is differentiable in $h$ with
\[
\varphi'(h) = f'(x_0 + h) - f'(x_0 - h).
\]

The numerator $\varphi(h)$ and denominator $h^2$ both vanish at $h = 0$, and $\frac{d}{dh}(h^2) = 2h \ne 0$ for $h \ne 0$. Apply L'H\^opital's rule once:
\[
\lim_{h \to 0} \frac{\varphi(h)}{h^2}
= \lim_{h \to 0} \frac{f'(x_0 + h) - f'(x_0 - h)}{2h}
= \frac{1}{2}\lim_{h \to 0}\!\left[\frac{f'(x_0+h) - f'(x_0)}{h} + \frac{f'(x_0-h) - f'(x_0)}{-h}\right]
\]
\[
= \frac{1}{2}\bigl[f''(x_0) + f''(x_0)\bigr] = f''(x_0),
\]
where the split into two difference quotients is valid because each individual limit exists (= $f''(x_0)$).

\[
\boxed{\lim_{h \to 0} \frac{1}{h^2}\left[f(x_0 - h) - 2f(x_0) + f(x_0 + h)\right] = f''(x_0).} \qquad \qed
\]


%% PROBLEM 2

\textbf{Problem 2.}

\[
I = \int_1^{1.5} x^2 \ln x\,dx, \qquad f(x) = x^2 \ln x.
\]

Derivatives:
\[
f''(x) = 2\ln x + 3, \qquad f^{(4)}(x) = -\frac{2}{x^2}.
\]

\textit{Exact value.} Integrating by parts ($u = \ln x$, $dv = x^2\,dx$):
\[
\int x^2 \ln x\,dx = \frac{x^3}{3}\ln x - \frac{x^3}{9} + C
\]
\[
I = \left[\frac{(1.5)^3}{3}\ln 1.5 - \frac{(1.5)^3}{9}\right] - \left[0 - \frac{1}{9}\right]
= \frac{9}{8}\ln\!\frac{3}{2} - \frac{19}{72} \approx 0.1922593577.
\]

\textit{Trapezoidal rule.} $h = b-a = 0.5$, $x_0 = 1$, $x_1 = 1.5$:
\[
T = \frac{h}{2}[f(x_0) + f(x_1)] = \frac{0.5}{2}\!\left[0 + \tfrac{9}{4}\ln\!\tfrac{3}{2}\right] = \tfrac{9}{16}\ln\!\tfrac{3}{2} \approx 0.2280741.
\]
Actual error: $|I - T| \approx 3.5815 \times 10^{-2}$.

Error bound (with $h = b-a$): $|E_T| \le \dfrac{h^3}{12}\max_{[1,1.5]}|f''|$. Since $f''$ is increasing,
$\max|f''| = f''(1.5) = 2\ln 1.5 + 3 \approx 3.8109$.
\[
|E_T| \le \frac{(0.5)^3}{12}(2\ln 1.5 + 3) \approx 3.9697 \times 10^{-2}.
\]

\textit{Simpson's rule.} $h = 0.25$, $x_0 = 1$, $x_1 = 1.25$, $x_2 = 1.5$:
\[
S = \frac{h}{3}[f(x_0) + 4f(x_1) + f(x_2)]
= \frac{1}{12}\!\left[0 + 4(1.5625)\ln 1.25 + 2.25\ln 1.5\right] \approx 0.1922453.
\]
Actual error: $|I - S| \approx 1.4050 \times 10^{-5}$.

Error bound: $|E_S| \le \dfrac{h^5}{90}\max_{[1,1.5]}|f^{(4)}|$. Since $|f^{(4)}(x)| = 2/x^2$ is decreasing,
$\max|f^{(4)}| = |f^{(4)}(1)| = 2$.
\[
|E_S| \le \frac{(0.25)^5}{90}(2) \approx 2.1701 \times 10^{-5}.
\]

\textit{Comparison.}
\begin{center}
\begin{tabular}{lccc}
\toprule
Rule & Approximation & Actual error & Error bound \\
\midrule
Trapezoidal & $0.22807412$ & $3.5815 \times 10^{-2}$ & $3.9697 \times 10^{-2}$ \\
Simpson's   & $0.19224531$ & $1.4050 \times 10^{-5}$ & $2.1701 \times 10^{-5}$ \\
\bottomrule
\end{tabular}
\end{center}

In both cases, $|\text{actual error}| \le |\text{error bound}|$, as required.


%% PROBLEM 3

\textbf{Problem 3.}

Let $x_1 = x_0 + h$, $x_2 = x_0 + 2h$. We seek $a_0, a_1, a_2, k$ such that
\[
\int_{x_0}^{x_2} f(x)\,dx = a_0 f(x_0) + a_1 f(x_1) + a_2 f(x_2) + k\,f^{(4)}(\xi).
\]

The coefficients $a_i$ are independent of $x_0$, so each test polynomial $f(x) = x^n$ must give an identity in $x_0$. Equate coefficients of powers of $x_0$.

\textit{$n = 1$:} $f(x) = x$, $f^{(4)} \equiv 0$.
\[
\int_{x_0}^{x_0+2h} x\,dx = 2 x_0 h + 2h^2,
\qquad
\sum a_i (x_0 + ih) = (a_0+a_1+a_2)\,x_0 + (a_1 + 2a_2)\,h.
\]
\[
a_0 + a_1 + a_2 = 2h, \qquad a_1 + 2a_2 = 2h. \tag{I, II}
\]

\textit{$n = 2$:} $f(x) = x^2$, $f^{(4)} \equiv 0$.
\[
\int_{x_0}^{x_0+2h} x^2\,dx = 2x_0^2 h + 4 x_0 h^2 + \tfrac{8}{3}h^3,
\]
\[
\sum a_i (x_0+ih)^2 = (a_0+a_1+a_2)\,x_0^2 + (2a_1 + 4a_2)\,x_0 h + (a_1 + 4a_2)\,h^2.
\]
The coefficients of $x_0^2$ and $x_0$ reproduce (I), (II); the constant term gives:
\[
a_1 + 4 a_2 = \tfrac{8}{3}h. \tag{III}
\]

\textit{$n = 3$:} expanding $\sum a_i (x_0+ih)^3$ produces only one new equation (the constant-in-$x_0$ coefficient): $a_1 + 8 a_2 = 4h$. With $a_1 = 4h/3$ and $a_2 = h/3$ obtained below, $4h/3 + 8h/3 = 4h$ holds automatically. So $n = 3$ contributes no new constraint, and Simpson's rule has degree of precision $\ge 3$.

\textit{Solving (I)--(III):} (III) $-$ (II) $\Rightarrow 2a_2 = 2h/3 \Rightarrow a_2 = h/3$;
then (II) $\Rightarrow a_1 = 2h - 2(h/3) = 4h/3$;
and (I) $\Rightarrow a_0 = 2h - 4h/3 - h/3 = h/3$.
\[
\boxed{a_0 = \frac{h}{3}, \qquad a_1 = \frac{4h}{3}, \qquad a_2 = \frac{h}{3}.}
\]

\textit{Finding $k$.} Apply to $f(x) = x^4$, for which $f^{(4)}(x) = 24$ (constant, so $f^{(4)}(\xi) = 24$). Take $x_0 = 0$ (translation invariance):
\[
\text{LHS} = \int_0^{2h} x^4\,dx = \frac{32 h^5}{5},
\]
\[
\text{RHS} = \frac{h}{3}\cdot 0 + \frac{4h}{3}\cdot h^4 + \frac{h}{3}\cdot 16 h^4 + 24 k = \frac{4h^5}{3} + \frac{16 h^5}{3} + 24 k = \frac{20 h^5}{3} + 24 k.
\]
\[
\frac{32 h^5}{5} = \frac{20 h^5}{3} + 24 k \;\Longrightarrow\; 24 k = \frac{96 h^5 - 100 h^5}{15} = -\frac{4 h^5}{15} \;\Longrightarrow\; \boxed{k = -\frac{h^5}{90}.}
\]

\textit{Simpson's rule with error term:}
\[
\boxed{\int_{x_0}^{x_2} f(x)\,dx = \frac{h}{3}\left[f(x_0) + 4 f(x_1) + f(x_2)\right] - \frac{h^5}{90}\,f^{(4)}(\xi).}
\]


%% PROBLEM 4

\textbf{Problem 4.}

\textit{Definition 4.1 (textbook, p.\ 195):} The degree of precision of a quadrature formula is the largest positive integer $n$ such that the formula is exact for $x^k$ for each $k = 0, 1, \ldots, n$.

Let $E(f) = \int_a^b f(x)\,dx - \sum_{i=0}^{m} a_i f(x_i)$ denote the error functional. Both the integral and the finite sum are linear in $f$, so
\[
E[\alpha f + \beta g] = \alpha E[f] + \beta E[g], \quad \forall\, \alpha, \beta \in \mathbb{R},
\]
i.e., $E$ is linear. By definition, ``the formula is exact for $f$'' $\iff E(f) = 0$.

\textit{Claim:} The formula has degree of precision $n$ if and only if $E(x^k) = 0$ for all $k = 0, 1, \ldots, n$ and $E(x^{n+1}) \ne 0$.

\textit{Proof.}

$(\Longrightarrow)$ Suppose the formula has degree of precision $n$. By Definition 4.1, the formula is exact for $x^k$ for each $k = 0, 1, \ldots, n$; that is,
\[
E(x^k) = 0, \qquad k = 0, 1, \ldots, n.
\]
Suppose for contradiction that $E(x^{n+1}) = 0$. Then for every polynomial $P(x) = \sum_{k=0}^{n+1} c_k x^k$ of degree $\le n+1$, linearity gives
\[
E(P) = \sum_{k=0}^{n+1} c_k\, E(x^k) = 0,
\]
so the formula is exact for every polynomial of degree $\le n+1$ — in particular for $x^{n+1}$. This contradicts the maximality of $n$. Hence $E(x^{n+1}) \ne 0$.

$(\Longleftarrow)$ Suppose $E(x^k) = 0$ for $k = 0, 1, \ldots, n$ and $E(x^{n+1}) \ne 0$.

For any polynomial $P$ of degree $\le n$, write $P(x) = \sum_{k=0}^{n} c_k x^k$. By linearity,
\[
E(P) = \sum_{k=0}^{n} c_k E(x^k) = 0,
\]
so the formula is exact for every polynomial of degree $\le n$.

The polynomial $x^{n+1}$ has degree $n+1$, and $E(x^{n+1}) \ne 0$ shows the formula is not exact for it. Therefore $n$ is the largest integer for which the formula is exact for $x^k$ for all $k = 0, 1, \ldots, n$, so the degree of precision is $n$. \qed


%% PROBLEM 5

\textbf{Problem 5.}

\[
I = \int_1^{2} x \ln x\,dx, \qquad f(x) = x \ln x.
\]
Derivatives: $f''(x) = 1/x$, $\;f^{(4)}(x) = 2/x^3$. On $[1,2]$:
\[
\max_{[1,2]} |f''(x)| = f''(1) = 1, \qquad \max_{[1,2]} |f^{(4)}(x)| = f^{(4)}(1) = 2.
\]
Set tolerance $\varepsilon = 10^{-5}$, $b - a = 1$.

\textbf{(a) Composite Trapezoidal rule} ($h = (b-a)/n$).
\[
|E_T| \le \frac{(b-a)\,h^2}{12}\,\max|f''| = \frac{1}{12 n^2}.
\]
Require $\dfrac{1}{12 n^2} \le 10^{-5}$, i.e., $n^2 \ge \dfrac{10^5}{12} \approx 8333.33$, so $n \ge 91.288$.
Check $n = 91$: $1/(12 \cdot 91^2) \approx 1.006 \times 10^{-5} > 10^{-5}$ (fails); $n = 92$: $1/(12 \cdot 92^2) \approx 9.85 \times 10^{-6} \le 10^{-5}$.
\[
\boxed{n = 92, \qquad h = \frac{1}{92} \approx 0.010870 \quad (|E_T| \le 9.85 \times 10^{-6}).}
\]

\textbf{(b) Composite Simpson's rule} ($n$ must be even since subintervals are paired; $h = (b-a)/n$).
\[
|E_S| \le \frac{(b-a)\,h^4}{180}\,\max|f^{(4)}| = \frac{2}{180\, n^4} = \frac{1}{90 n^4}.
\]
Require $\dfrac{1}{90 n^4} \le 10^{-5}$, i.e., $n^4 \ge \dfrac{10^5}{90} = \dfrac{10^4}{9} \approx 1111.11$, so $n \ge 5.774$.
Check $n = 4$ (even): $1/(90 \cdot 256) \approx 4.34 \times 10^{-5} > 10^{-5}$ (fails); $n = 6$: $1/(90 \cdot 1296) \approx 8.57 \times 10^{-6} \le 10^{-5}$.
\[
\boxed{n = 6, \qquad h = \frac{1}{6} \approx 0.166667 \quad (|E_S| \le 8.57 \times 10^{-6}).}
\]

\textbf{(c) Composite Midpoint rule} ($n$ even since midpoints come in paired blocks; $h = (b-a)/(n+2)$, B\&F 10th ed.\ Eq.~(4.34)).
\[
|E_M| \le \frac{(b-a)\,h^2}{6}\,\max|f''| = \frac{1}{6 (n+2)^2}.
\]
Require $\dfrac{1}{6 (n+2)^2} \le 10^{-5}$, i.e., $(n+2)^2 \ge \dfrac{10^5}{6} \approx 16666.67$, so $n + 2 \ge 129.10$, i.e., $n \ge 127.10$.
Check $n = 126$ (even, $n+2=128$): $1/(6 \cdot 128^2) \approx 1.017 \times 10^{-5} > 10^{-5}$ (fails); $n = 128$ ($n+2 = 130$): $1/(6 \cdot 130^2) \approx 9.86 \times 10^{-6} \le 10^{-5}$.
\[
\boxed{n = 128, \qquad h = \frac{1}{130} \approx 0.007692 \quad (|E_M| \le 9.86 \times 10^{-6}).}
\]


%% PROBLEM 6

\textbf{Problem 6.}

With $h_1 = b - a$ (so $n = 1$) and $h_2 = (b-a)/2$ (so $n = 2$), the composite Trapezoidal rule (1) gives:
\begin{align*}
R_{1,1} &= \frac{b - a}{2}\,[f(a) + f(b)], \\[3pt]
R_{2,1} &= \frac{b - a}{4}\,\!\left[f(a) + 2 f\!\left(\tfrac{a+b}{2}\right) + f(b)\right].
\end{align*}

Richardson's extrapolation:
\[
R_{2,2} = R_{2,1} + \frac{1}{3}(R_{2,1} - R_{1,1}).
\]

Let $A = f(a)$, $M = f((a+b)/2)$, $B = f(b)$. Then:
\begin{align*}
R_{2,1} - R_{1,1} &= \frac{b-a}{4}(A + 2M + B) - \frac{b-a}{2}(A + B)
= \frac{b-a}{4}(-A + 2M - B), \\[3pt]
\frac{1}{3}(R_{2,1} - R_{1,1}) &= \frac{b-a}{12}(-A + 2M - B), \\[3pt]
R_{2,2} &= \frac{b-a}{4}(A + 2M + B) + \frac{b-a}{12}(-A + 2M - B) \\
       &= \frac{b-a}{12}\!\left[3(A + 2M + B) + (-A + 2M - B)\right] \\
       &= \frac{b-a}{12}\,(2A + 8M + 2B) = \frac{b - a}{6}\,(A + 4M + B).
\end{align*}

\[
\boxed{R_{2,2} = \frac{b-a}{6}\!\left[f(a) + 4 f\!\left(\tfrac{a+b}{2}\right) + f(b)\right].}
\]

With $h = (b-a)/2$, this is exactly Simpson's rule:
\[
R_{2,2} = \frac{h}{3}\!\left[f(a) + 4 f\!\left(\tfrac{a+b}{2}\right) + f(b)\right].
\]


%% PROBLEM 7

\textbf{Problem 7.}

The IVP is $y' = [\exp(t)\,y]^2 = e^{2t} y^2$, $\;y(0) = -1$, $\;0 \le t \le 1$, with
\[
f(t, y) = e^{2t}\, y^2, \qquad D = \{(t, y) : 0 \le t \le 1, \; -\infty < y < \infty\}.
\]

\textit{Continuity.} $f$ is a polynomial in $y$ multiplied by a continuous function of $t$, hence continuous on $D$. \checkmark

\textit{Theorem 5.3 cannot establish Lipschitz.} Compute $\partial f / \partial y = 2 e^{2t} y$. On $D$ the variable $y$ ranges over all of $\mathbb{R}$, so
\[
\sup_{(t,y) \in D} \left|\frac{\partial f}{\partial y}\right| = \sup_{0 \le t \le 1} \sup_{y \in \mathbb{R}} 2 e^{2t} |y| = +\infty.
\]
Theorem 5.3 (which gives a Lipschitz constant $L = \sup_D |\partial f/\partial y|$) requires this supremum to be finite, so it does not yield a Lipschitz constant for $f$ on $D$. But Theorem 5.3 is only sufficient; we must check the Lipschitz condition directly.

\textit{Direct verification that $f$ is not Lipschitz on $D$.} Suppose, for contradiction, that there exists $L > 0$ with $|f(t, y_1) - f(t, y_2)| \le L\,|y_1 - y_2|$ for all $(t, y_1), (t, y_2) \in D$. Take $t = 0$, $y_2 = 0$, and any $y_1 = M > 0$:
\[
\frac{|f(0, M) - f(0, 0)|}{|M - 0|} = \frac{|M^2 - 0|}{M} = M.
\]
This must be $\le L$, but $M$ is unrestricted in $D$. Contradiction. Hence $f$ does \textbf{not} satisfy a Lipschitz condition on $D$, so the hypothesis of Theorem 5.4 fails.

\[
\boxed{\text{Theorem 5.4 \textbf{cannot} be applied to (2) on } D = \{0 \le t \le 1,\ y \in \mathbb{R}\}.}
\]

\textit{Remark.} A unique solution $y(t) = -2/(1 + e^{2t})$ does exist (Theorem 5.4 is only sufficient). Verification:
\[
y'(t) = -2 \cdot \frac{-2 e^{2t}}{(1 + e^{2t})^2} = \frac{4 e^{2t}}{(1 + e^{2t})^2} = e^{2t}\!\left(\frac{-2}{1 + e^{2t}}\right)^{\!2} = e^{2t} y^2 = [\exp(t)\, y]^2,
\]
and $y(0) = -2/2 = -1$. Since $y(0) = -1$ and $y$ increases monotonically toward $0$ on $[0,1]$, $|y(t)| \le 1$ throughout. Restricting $D$ to a strip $\{|y| \le M\}$ would render $\partial f/\partial y$ bounded, and Theorem 5.4 would then apply.


%% PROBLEM 8

\textbf{Problem 8.}

$y' = t(y + 1)$, $\;y(0) = 0$, $\;0 \le t \le 1$, $\;h = 1/3$.

Euler's method: $w_{i+1} = w_i + h \cdot f(t_i, w_i) = w_i + h\, t_i (w_i + 1)$, with $t_i = i h$.

\begin{align*}
w_0 &= 0, &t_0 &= 0, \\[3pt]
w_1 &= w_0 + \tfrac{1}{3} \cdot 0 \cdot (0 + 1) = 0, &t_1 &= \tfrac{1}{3}, \\[3pt]
w_2 &= w_1 + \tfrac{1}{3} \cdot \tfrac{1}{3} \cdot (0 + 1) = \tfrac{1}{9}, &t_2 &= \tfrac{2}{3}, \\[3pt]
w_3 &= w_2 + \tfrac{1}{3} \cdot \tfrac{2}{3} \cdot \!\left(\tfrac{1}{9} + 1\right)
     = \tfrac{1}{9} + \tfrac{2}{9} \cdot \tfrac{10}{9}
     = \tfrac{1}{9} + \tfrac{20}{81}
     = \tfrac{9 + 20}{81}
     = \tfrac{29}{81}, &t_3 &= 1.
\end{align*}

\begin{center}
\begin{tabular}{cccc}
\toprule
$i$ & $t_i$ & $w_i$ & $y(t_i) = e^{t_i^2/2} - 1$ \\
\midrule
0 & $0$       & $0$           & $0$ \\
1 & $1/3$     & $0$           & $e^{1/18} - 1 \approx 0.05709$ \\
2 & $2/3$     & $1/9$         & $e^{2/9} - 1 \approx 0.24860$ \\
3 & $1$       & $29/81$       & $e^{1/2} - 1 \approx 0.64872$ \\
\bottomrule
\end{tabular}
\end{center}

\[
\boxed{w_0 = 0, \quad w_1 = 0, \quad w_2 = \tfrac{1}{9}, \quad w_3 = \tfrac{29}{81}.}
\]

\par
\end{document}


%%% Local Variables:
%%% mode: latex
%%% TeX-master: t
%%% End:
